% ==============================================================================
% Section 1: Introduction
% ==============================================================================

\section{Introduction}
\label{sec:introduction}

\subsection{The collapse of a keystone predator}

Between 2013 and 2015, the largest marine wildlife mass mortality event ever recorded
swept the Pacific coast of North America. Sea star wasting disease (SSWD) killed
billions of sea stars across more than twenty species, from Alaska to Baja California.
No species was hit harder than the sunflower star, \textit{Pycnopodia helianthoides}---the
largest sea star on Earth, capable of spanning a full meter across, with up to
twenty-four arms and thousands of tube feet that let it glide across the seafloor at
speeds remarkable for an echinoderm. Survey data compiled by \citet{gravem2021}
documented a 90.6\% range-wide decline, a loss so severe that \textit{Pycnopodia} was
subsequently listed as Critically Endangered on the IUCN Red List---the first sea star
to receive that designation.

The ecological consequences of this collapse extend far beyond the loss of a single
species. \textit{Pycnopodia} is a keystone predator in kelp forest ecosystems, where it
controls populations of sea urchins---the primary grazers of kelp. Without sunflower
stars to keep urchin numbers in check, urchin populations have exploded across large
stretches of the coast, overgrazing kelp forests and converting them into ``urchin
barrens'': rocky reefs stripped of macroalgae, impoverished in biodiversity, and
largely devoid of the three-dimensional habitat structure that supports hundreds of
associated species. The loss of \textit{Pycnopodia} has thus triggered a trophic cascade
with consequences that ripple through entire nearshore ecosystems, from fish nursery
habitat to carbon sequestration.

\subsection{Identifying the pathogen}

For over a decade after the initial outbreak, the causative agent of SSWD remained
uncertain. Early hypotheses centered on a densovirus (sea star--associated
densovirus, SSaDV), but subsequent work failed to establish a consistent causal
link. The breakthrough came when \citet{prentice2025}, publishing in \textit{Nature
Ecology \& Evolution}, fulfilled Koch's postulates for the marine bacterium
\textit{Vibrio pectenicida}. Their work demonstrated that this Vibrio species could
be isolated from wasting sea stars, grown in pure culture, used to experimentally
induce wasting in healthy animals, and then re-isolated from the newly diseased
individuals---the classical chain of evidence required to establish causation in
infectious disease.

The identification of \textit{V.\ pectenicida} as the etiological agent transforms our
understanding of SSWD from a mysterious syndrome into a problem that can, in
principle, be modeled mechanistically. Vibrio species are well-studied bacteria
whose ecology is tightly coupled to seawater temperature: they enter a viable but
non-culturable (VBNC) dormant state when water temperatures drop below a
species-specific thermal threshold, and resume active growth and virulence when
temperatures rise above it. This temperature-dependent life cycle means that climate
warming does not merely provide a stressful background for the host---it directly
controls the seasonal window during which the pathogen can cause disease.

\subsection{The conservation challenge}

Active conservation efforts for \textit{Pycnopodia} are already underway. Captive
breeding programs, informed by detailed work on sunflower star spawning biology and
larval rearing \citep{hodin2021}, aim to produce juveniles for eventual
reintroduction to depleted habitats. Field researchers have identified potential
refugia---\citet{gehman2025} documented that fjord environments in British Columbia,
where a surface freshwater lens suppresses \textit{Vibrio} activity, may harbor
surviving populations shielded from the worst disease pressure.

Yet these efforts face a fundamental uncertainty: \textit{what will the disease
landscape look like in the decades ahead?} Captive-bred juveniles released into the
wild will encounter a pathogen that is itself evolving---potentially adapting to new
thermal regimes, shifting its virulence, and tracking the warming ocean northward.
Managers selecting reintroduction sites need to know not just where \textit{Pycnopodia}
survives today, but where it can persist through mid-century under realistic climate
and evolutionary trajectories. Without forecasts that account for the coupled dynamics
of host populations, pathogen evolution, and climate change, conservation investments
risk targeting sites that will become untenable within a generation.

\subsection{Why an eco-evolutionary model?}

Classical epidemiological models treat host susceptibility and pathogen virulence as
fixed parameters. For a disease system operating over decades in a warming ocean,
this assumption is untenable. Host populations under intense disease pressure will
experience natural selection for resistance---individuals that survive carry alleles
conferring some degree of tolerance or resistance, and pass them to their offspring.
Simultaneously, the pathogen population evolves: \textit{V.\ pectenicida} can adapt
its thermal tolerance (the VBNC threshold temperature) to track warming waters, and
its virulence can shift in response to the trade-off between transmission success
and host mortality \citep{anderson1982}. These evolutionary dynamics operate on
timescales of years to decades---precisely the timescale relevant to conservation
planning.

We therefore developed a coupled eco-evolutionary epidemiological agent-based model
(ABM) that simulates individual \textit{Pycnopodia} across 896 discrete coastal sites
spanning the species' entire range, from the Aleutian Islands of Alaska to Baja
California, Mexico. Each simulated sea star carries a heritable resistance trait;
each local pathogen population carries heritable traits for thermal tolerance and
virulence. Reproduction, mortality, disease transmission, and trait inheritance all
operate at the individual level, allowing emergent population dynamics and
evolutionary trajectories to arise from first principles rather than imposed
aggregate equations.

\subsection{Scope of this report}

We calibrated the model against observational survey data from 2012 through 2025,
achieving a root-mean-square error (RMSE) of 0.504 against normalized regional
abundance indices. We then projected the model forward to 2050 under four scenarios
designed to isolate the effects of climate trajectory and pathogen evolutionary
capacity:

\begin{description}
    \item[Scenario F01 (Baseline):] Moderate warming (Shared Socioeconomic Pathway 2-4.5, hereafter SSP2-4.5) with all pathogen
    evolutionary mechanisms active---thermal adaptation and virulence evolution.
    This represents our best estimate of the most likely future trajectory.

    \item[Scenario F02 (No pathogen evolution):] Moderate warming (Shared Socioeconomic Pathway 2-4.5, hereafter SSP2-4.5) with
    all pathogen evolution disabled. The pathogen retains its initial thermal
    tolerance and virulence throughout the simulation. Comparing F02 to F01 reveals
    the total effect of pathogen evolutionary capacity on host population outcomes.

    \item[Scenario F03 (Thermal adaptation only):] Moderate warming (Shared Socioeconomic Pathway 2-4.5, hereafter SSP2-4.5)
    with pathogen thermal adaptation enabled but virulence evolution disabled.
    Comparing F03 to F01 isolates the specific contribution of virulence evolution.

    \item[Scenario F04 (High warming):] High-end warming (SSP5-8.5) with all
    pathogen evolution active. Comparing F04 to F01 reveals how an accelerated
    warming trajectory reshapes outcomes across the species' latitudinal range.
\end{description}

Each scenario was run with three independent random seeds, and all results reported
here represent the three-seed mean. The peak historical population across all
scenarios was 4,447,241 individuals.

This report presents the key findings from these forecast simulations (Section~\ref{sec:key_findings}), describes the scenario design (Section~\ref{sec:scenarios}), provides detailed regional results (Section~\ref{sec:results}), explains the model architecture (Section~\ref{sec:model}), and discusses implications for conservation strategy (Section~\ref{sec:discussion}). Our aim is to provide the quantitative foundation
that conservation managers, captive breeding programs, and policy makers need to
make informed decisions about where, when, and how to invest in the recovery of
this critically endangered keystone predator.
